\documentclass[../../../uem_phy_std_q.tex]{subfiles}

\begin{document}
    金属導線中の電子の運動について考える．断面積$S$，長さ$L$の円筒状の導線の%
    両端A，B間に電位差（電圧）$V$を加える．%
    ただし，Aの電位よりBの電位の方が高いものとする．%
    また，電子の質量を$m$，電子の電荷を$-e$，導線の単位体積あたりの自由電子の数を$n$とする．%
    自由電子は電場の存在下では平均するとAB方向に運動しているので，%
    以下では1次元の運動として考える．また，導線中を移動する自由電子には%
    速さに比例する抵抗力が働くものとし，比例定数を$k$とする．
    
    \begin{enumerate}
        \item%
            自由電子の運動の向きはA→BとB→Aのどちらであるか答えよ．
        \item%
            自由電子が導線中を平均して一定の速さ$v$で動いているとする．%
            $v$を求めよ．
        \item%
            導線を流れる電流の大きさ$I$を求めよ．
        \item%
            この導線の抵抗率$\rho$を求めよ．
        \item%
            この導線内の全自由電子が電場からされる単位時間あたりの仕事（仕事率）$P$を，%
            $I$，$V$を用いて表せ．
    \end{enumerate}
\end{document}



